\chapter{Atlas Stream Processing and Event-Driven Pipelines}
\index{Atlas Stream Processing}\index{stream processing}\index{tumbling window}\index{hopping window}\index{session window}\index{checkpointing}\index{dead-letter queue}\index{exactly-once}\index{stream join}%

\begin{objectives}
\begin{itemize}
    \item Explain Atlas Stream Processing (ASP) concepts: stream processors, workspaces, connections, and pipelines.
    \item Consume from Kafka, Kinesis, and MongoDB change streams inside ASP pipelines.
    \item Implement stream-collection joins with \texttt{\$lookup} and reason about time-windowed stream-stream joins and late arrivals.
    \item Choose tumbling, hopping, and session windows and manage window/checkpoint state growth.
    \item Build end-to-end exactly-once-effect sinks with checkpointing plus idempotent \texttt{\$merge} upserts on a dedup key, with a dead-letter queue for poison events.
    \item Architect a real-time fraud-detection pipeline in the Friday engagement.
\end{itemize}
\end{objectives}

\begin{crosscloud}
Managed streaming with SQL-like or pipeline semantics over a broker backbone is a commodity layer now; the differentiators are state management, exactly-once composition, and how tightly the sink integrates with the operational store.

\vspace{2mm}
{\footnotesize
\begin{tabular}{@{}p{2.9cm}p{3.0cm}p{3.0cm}p{3.0cm}@{}} 
\toprule
\textbf{Capability} & \textbf{AWS} & \textbf{Azure} & \textbf{GCP} \\
\midrule
Broker backbone & MSK / Kinesis & Event Hubs (Kafka API) & Pub/Sub / Managed Kafka \\
Stream processor & Kinesis Data Analytics (Flink) & Stream Analytics / Fabric RTI & Dataflow (Beam) \\
Windowing & Tumbling/hopping/session & Tumbling/hopping/session & Full Beam window model \\
State/checkpointing & Flink state backends & Managed & Managed (streaming engine) \\
Exactly-once sink & Idempotent/transactional sinks & Idempotent upserts & Idempotent writes \\
DB-native option & DMS streams & Cosmos change feed + Functions & Firestore triggers + Dataflow \\
\addlinespace
Snowflake / Fabric & \multicolumn{3}{l}{Streams + Tasks for incremental CDC pipelines; Snowpipe Streaming; Fabric eventstreams with KQL windows} \\
\bottomrule
\end{tabular}}

\vspace{1mm}
ASP's distinction is the same as Chapter 13's: the streaming tier speaks the aggregation pipeline and sinks directly into Atlas collections, so enrichment joins and materialized outputs stay on one platform \cite{mongodbStreamProcessingDocs,apacheKafkaRepo}.
\end{crosscloud}

\section{Week 15 Roadmap: Continuous Computation Over the Firehose}
Chapter 16 will consume change streams imperatively; this week computes over streams declaratively. Atlas Stream Processing runs continuous aggregation pipelines over event sources---Kafka topics, Kinesis streams, or MongoDB change streams---with windows, joins, state, and checkpointing managed for you. The week's spine is the syllabus's hardest requirement: making the sink end-to-end exactly-once in effect, because at-least-once delivery is the honest model of every broker \cite{mongodbStreamProcessingDocs,kleppmann2017ddia}.

\begin{definitionbox}
A \textbf{stream processor} in ASP is a continuously running aggregation pipeline with a \texttt{\$source} (a connection to a broker or change stream), optional windowing and enrichment stages, and a sink (\texttt{\$merge} or \texttt{\$emit}), checkpointed by the platform for fault tolerance.
\end{definitionbox}

\begin{keyterms}
\textbf{Connection}: a named, reusable source/sink binding (Kafka cluster, Atlas database). \textbf{Window}: tumbling (fixed, non-overlapping), hopping (fixed, overlapping), or session (gap-closed) event-time grouping. \textbf{Checkpoint}: persisted processor state enabling resume-after-failure. \textbf{State store}: the per-key window state whose growth must be bounded. \textbf{DLQ}: dead-letter queue for poison events. \textbf{Exactly-once effect}: observable state equivalent to single application, achieved by idempotent sinks.
\end{keyterms}

\section{Monday: ASP Concepts}
\subsection{Workspaces, Connections, and Processors}
An ASP workspace hosts connections (Kafka clusters, Atlas databases, HTTPS endpoints) and stream processors. A processor's pipeline is the aggregation language you already know---\texttt{\$match}, \texttt{\$set}, \texttt{\$lookup}, \texttt{\$group} inside windows---extended with streaming primitives: \texttt{\$source} reads a connection, \texttt{\$tumblingWindow}/\texttt{\$hoppingWindow} group by event time, \texttt{\$merge} upserts into an Atlas collection, and \texttt{\$emit} writes to a topic. Validation is a first-class operation: \texttt{sp.process()} starts the processor, and a ``sample'' mode runs the pipeline against live source data without committing a sink---always sample before starting.

\subsection{The Operational Contract}
ASP manages checkpoints automatically: processor state (window contents, source offsets) is persisted at intervals, and a restarted processor resumes from its last checkpoint. The interval is a latency/recovery trade-off: frequent checkpoints bound duplicate re-emission after a crash but cost throughput; infrequent checkpoints do the opposite. What ASP does \emph{not} do is make your sink idempotent---that remains your schema design (Wednesday).

\begin{notebox}
A change-stream source gives ASP the database's own event feed (Chapter 16 mechanics: resume tokens, oplog window). The same sizing rule applies: if the processor pauses longer than the oplog window, the stream cannot resume and must re-bootstrap from a snapshot---size the oplog for your worst-case processor downtime.
\end{notebox}

\section{Tuesday: Streaming Ingestion and Stream Joins}
\subsection{Sources: Kafka, Kinesis, Change Streams}
Connection registry entries carry bootstrap servers, topics, authentication (SASL, IAM, x.509), and schema/serialization settings. The three source families cover the platform patterns: Kafka topics for application event buses, Kinesis for AWS-native pipelines, and change streams for database-as-event-source designs. Source configuration includes the starting position (earliest vs.\ latest)---a deliberate choice: reprocessing history is sometimes the feature (backfill) and sometimes the bug (double-sending notifications).

\subsection{Stream-Collection Enrichment}
\texttt{\$lookup} inside an ASP pipeline joins each event against an Atlas collection---the fraud-pipeline pattern: enrich a card event with cardholder data before scoring. The join is a point lookup per event against the current collection state; it is not a temporal join (the event sees the cardholder as of now, not as of event time). For correctness under updates, either accept current-state semantics (usually right for enrichment) or model the dimension as an event stream and do a windowed stream-stream join.

\begin{lstlisting}[language=JavaScript, caption={ASP pipeline skeleton: source, enrich, dedup-key, idempotent sink.}]
pipeline = [
  { $source: { connectionName: "kafka_orders" } },
  { $lookup: {                          // stream-collection enrichment
      from: { connectionName: "atlas_db", db: "shop", coll: "customers" },
      localField: "customer_id", foreignField: "_id", as: "customer" } },
  { $set: { dedup_key: { $concat: ["$customer_id", ":", "$event_id"] } } },
  { $merge: {                           // exactly-once-effect sink
      into: { connectionName: "atlas_db", db: "shop", coll: "enriched_orders" },
      on: "dedup_key",
      whenMatched: "replace", whenNotMatched: "insert" } }
];
\end{lstlisting}

\subsection{Stream-Stream Joins and Late Arrivals}
Joining two event streams requires a time window: an event waits for its counterpart within the window or matches nothing. Late arrivals---events whose event time precedes the window's watermark---are the defining complication: within the allowed lateness they join and correct the output; beyond it they are dropped or routed to a side output. The design questions are the watermark (how late is ``late''), the join window (how long events wait), and the correction policy (update the emitted result, or emit a compensating event).

\section{Wednesday: Windowing and State Management}
\subsection{Tumbling, Hopping, and Session Windows}
\textbf{Tumbling} windows partition time into fixed non-overlapping buckets (five-minute sales totals). \textbf{Hopping} windows overlap (a five-minute window every minute) for smoother moving metrics. \textbf{Session} windows close after an inactivity gap, grouping a user's burst of activity. Event-time processing (timestamps from the events, not arrival time) is what makes results correct under out-of-order delivery; every window definition names its time field and its allowed lateness.

\subsection{State Growth Is the Killer}
Every open window holds state per group key. High-cardinality keys---one window per user per five minutes across ten million users---grow the state store unboundedly, and the syllabus's warning is literal: unbounded state kills streaming jobs through memory pressure, slow checkpoints, and restart storms (each restart reloads and rebuilds state). The defenses: bound cardinality by design (windowed aggregates with watermark expiry, not per-key-forever state), monitor state size and checkpoint duration as first-class metrics, and treat a growing-state trend as a capacity incident.

\begin{pitfall}
\textbf{Windowed \texttt{\$merge} without a deterministic key.} At-least-once delivery means window emissions can repeat after a checkpoint restore. A sink keyed on \texttt{user\_id + window\_start} upserts idempotently; a sink without it duplicates rows on every replay. The deterministic \texttt{\_id} (window start + group key) is the entire exactly-once story.
\end{pitfall}

\subsection{Composing Exactly-Once Effect}
End-to-end exactly-once is a composition, not a feature flag: checkpointing gives the processor a consistent resume point; the source replays from it (at-least-once); the sink's idempotent upsert makes replay harmless; and the DLQ isolates poison events so one malformed message cannot stall the stream forever. Alert on consumer lag---not on errors---as the paging signal; a processor silently falling behind is the failure mode that page-worthy systems catch late.

\section{Thursday Hands-On Project: Exactly-Once Kafka Windowed Pipeline}
Build the syllabus lab: Kafka source, five-minute tumbling windows, idempotent \texttt{\$merge} into Atlas, DLQ for poison events.

\subsection{Lab Protocol}
\begin{enumerate}
    \item Stand up a local Kafka (or Redpanda) container and an Atlas/Enterprise connection registry in ASP.
    \item Produce a scripted order-event stream with two deliberate poison messages (malformed JSON, missing key field).
    \item Define the processor: source $\to$ parse/validate (route failures to a dead-letter topic) $\to$ five-minute tumbling window summing revenue per SKU $\to$ \texttt{\$merge} with \texttt{on: ["sku", "window\_start"]}.
    \item Prove exactly-once effect: kill the processor mid-window, restart it, and verify the output collection matches a single-run reference exactly (no duplicates, no gaps).
    \item Prove the DLQ: both poison events land in the dead-letter topic and the stream continues.
\end{enumerate}

\subsection*{Deliverables}
\begin{itemize}
    \item The connection and processor definitions as code, plus the producer script.
    \item The kill-and-restart evidence: output diff against the reference run.
    \item The DLQ contents and the lag/checkpoint metrics captured during the drill.
\end{itemize}

\subsection*{Acceptance Criteria}
\begin{itemize}
    \item Post-restart output is byte-identical to the uninterrupted reference run.
    \item Both poison events reach the DLQ with the stream uninterrupted.
    \item Checkpoint interval, state size, and lag are captured as metrics, not impressions.
\end{itemize}

\subsection*{Stretch Goals}
\begin{itemize}
    \item Switch the window to hopping (5-minute size, 1-minute hop) and quantify the state-size increase.
    \item Let the processor pause past a shrunken oplog window on a change-stream source and execute the re-bootstrap procedure (snapshot + fresh stream).
\end{itemize}

\section{Friday Use Case and Architecture: Real-Time Fraud Detection}
\subsection{Scenario and Requirements}
A card network must score authorization events in under 300 ms end-to-end: enrich each event with cardholder data, compute velocity features (authorizations per card per five minutes), and persist materialized alerts---with malformed events isolated and no duplicate alerts after infrastructure failures.

\subsection{Architecture}
The ASP pipeline: Kafka authorization events $\to$ validation with DLQ routing $\to$ \texttt{\$lookup} enrichment against the cardholders collection $\to$ a tumbling velocity window per card $\to$ a scoring stage that joins the velocity stream against per-card risk profiles (windowed stream-stream join with a two-minute lateness allowance) $\to$ \texttt{\$merge} alerts keyed on \texttt{card\_id + window\_start}. Exactly-once effect comes from checkpointing plus the deterministic merge key; cardholder freshness is a change-stream-fed cache inside Atlas so the enrichment lookup never leaves the platform. SLOs are wired per Chapter 23: end-to-end latency histogram, consumer lag, state size, checkpoint duration, and DLQ depth, each with a runbook entry.

\begin{figure}[H]
    \centering
    \resizebox{0.98\textwidth}{!}{%
    \begin{tikzpicture}[
        st/.style={stage, minimum width=2.15cm, font=\scriptsize\bfseries},
        bad/.style={rectangle, rounded corners, draw=red!60!black, fill=red!6, minimum width=2.0cm, minimum height=0.9cm, align=center, font=\scriptsize\bfseries},
        store/.style={cylinder, shape aspect=0.25, draw=primary, thick, fill=primary!8, shape border rotate=90, align=center, minimum width=1.5cm, minimum height=1.1cm, font=\scriptsize\bfseries},
        arrow/.style={-{Stealth[length=2.2mm]}, draw=primary, thick}
    ]
        \node[st] (src) at (0,0) {\$source\\Kafka};
        \node[st] (val) at (2.7,0) {validate\\+ route};
        \node[st] (enr) at (5.4,0) {\$lookup\\cardholder};
        \node[st] (win) at (8.1,0) {tumbling\\velocity 5m};
        \node[st] (score) at (10.8,0) {score +\\risk join};
        \node[store] (sink) at (13.6,0) {alerts\\(\$merge)};
        \node[bad] (dlq) at (2.7,-1.9) {DLQ topic};
        \draw[arrow] (src) -- (val);
        \draw[arrow] (val) -- (enr);
        \draw[arrow] (enr) -- (win);
        \draw[arrow] (win) -- (score);
        \draw[arrow] (score) -- (sink);
        \draw[arrow, dashed] (val) -- (dlq);
    \end{tikzpicture}%
    }
    \caption{Fraud pipeline: validate, enrich, window, score, and idempotently merge---with poison events routed to the DLQ.}
    \label{fig:ch15-fraud}
\end{figure}

\subsection{Guardrails}
Four standing controls: state-size and checkpoint-duration alerts (state growth is the slow killer), a lag-based page (not error-based), a monthly kill-and-restart drill re-proving the exactly-once effect, and a documented re-bootstrap procedure for oplog-window expiry on any change-stream source.

\section{Interview Q\&A}
\subsection{Streaming}
\begin{enumerate}
    \item \textbf{Q: Why does unbounded state from high-cardinality keys kill a streaming job?}\\
    A: Every open key holds memory and must be checkpointed; growth without bounds means memory pressure, slow checkpoints, and restarts that reload ever-larger state---a storm that never converges.

    \item \textbf{Q: What does \texttt{\$lookup} enable inside an ASP pipeline?}\\
    A: Stream-collection enrichment: each event joins current Atlas collection state (cardholders, product catalogs) without leaving the platform. It is current-state, not temporal---event-time-correct dimensions need stream-stream joins.

    \item \textbf{Q: What do late-arriving events imply for windowed joins?}\\
    A: Within allowed lateness they correct the window's output; beyond it they are dropped or side-output. Watermarks and join windows are business decisions about how late data may be.

    \item \textbf{Q: How do checkpointing + idempotent \texttt{\$merge} deliver exactly-once effect?}\\
    A: Checkpoints resume the processor consistently; replay re-emits windows (at-least-once); the sink upsert on a deterministic key (window start + group key) makes re-emission harmless.

    \item \textbf{Q: What happens when a resume token outlives the oplog window?}\\
    A: The stream cannot resume---the entries it needs have rolled off. Re-bootstrap: snapshot the source, open a fresh stream from the post-snapshot point, and resume checkpointing. Size the oplog beyond worst-case processor downtime.
\end{enumerate}

\section{Further Reading}
The Atlas Stream Processing documentation covers connections, processors, windows, and checkpointing \cite{mongodbStreamProcessingDocs}; Kafka's semantics underpin the source model \cite{apacheKafkaRepo}. Kleppmann's stream-processing chapters are the theory base \cite{kleppmann2017ddia}, and Chapter 16's change-stream consumer mechanics complete the picture (resume tokens, pre-images, imperative consumers) \cite{mongodbChangeStreamsDocs}.

\section*{Key Takeaways}
\begin{itemize}
    \item ASP runs continuous aggregation pipelines over Kafka, Kinesis, or change streams, with managed checkpoints.
    \item Sample before you start; validate before you merge; route poison events to a DLQ.
    \item Windows group by event time; watermarks and allowed lateness are business decisions.
    \item State growth is the slow killer: bound cardinality, monitor state and checkpoint duration.
    \item Exactly-once is a composition: checkpoint + at-least-once replay + deterministic-key upsert.
    \item Lag is the paging metric; a silently behind processor is worse than a crashing one.
\end{itemize}

\section{Exercises}
\begin{enumerate}
    \item \textbf{(Conceptual)} Choose tumbling, hopping, or session windows for: hourly revenue, a live dashboard's moving five-minute error rate, and per-user activity bursts.
    \item \textbf{(Conceptual)} Explain why \texttt{\$lookup} enrichment is not event-time-correct, and when that matters.
    \item \textbf{(Conceptual)} Design the dedup key for a sensor-rollup pipeline; show why each component is necessary.
    \item \textbf{(Hands-on)} Run the Thursday lab and capture the kill-and-restart parity evidence.
    \item \textbf{(Hands-on)} Measure state size and checkpoint duration as you raise group-key cardinality 10$\times$; plot the trend.
    \item \textbf{(Design)} Design the watermark and lateness policy for the fraud pipeline, with the business justification.
    \item \textbf{(Design)} Write the re-bootstrap runbook for a change-stream-sourced processor after oplog expiry.
    \item \textbf{(Interview)} ``Alerts doubled after every deploy until we fixed the sink.'' Explain the mechanism.
\end{enumerate}

\section{Capstone Project: Provably Idempotent Streaming Platform}\label{sec:ch15-capstone}

\begin{capstonebox}
\textbf{Mission:} deliver the fraud-pipeline pattern as a reusable platform: ASP processor with enrichment, windowing, exactly-once-effect sink, DLQ, and full SLO wiring---proven by a kill-and-restart drill and a poison-flood drill, with state-growth monitoring as a first-class alert.

\textbf{Estimated time:} 6--8 hours.

\textbf{What you will practice:}
\begin{itemize}
    \item Streaming pipeline design with enrichment and windowed joins.
    \item Exactly-once effect via checkpointing and deterministic merge keys.
    \item Operationalizing a stream processor: lag, state, checkpoints, DLQ.
\end{itemize}
\end{capstonebox}

\subsection*{Scenario}
The fraud platform's reviewers require evidence that alerts are never duplicated and never lost, including across processor failures and a flood of malformed events. This capstone is that evidence.

\subsection*{Requirements}
\begin{itemize}
    \item The full pipeline (source, validation+DLQ, enrichment, velocity window, risk join, \texttt{\$merge} sink) defined as code.
    \item A kill-and-restart drill proving output parity with an uninterrupted reference run.
    \item A poison-flood drill (100 malformed events) proving DLQ isolation and stream continuity.
    \item Metrics and alerts: consumer lag, state size, checkpoint duration, DLQ depth, end-to-end latency.
    \item A state-growth capacity note: measured state per key, projected growth, and the watermark-expiry bound.
\end{itemize}

\subsection*{Implementation Steps}
\begin{enumerate}
    \item \textbf{Build.} Define connections and processor; sample against live data before starting.
    \item \textbf{Prove idempotency.} Run the kill-and-restart drill; diff outputs.
    \item \textbf{Prove isolation.} Run the poison flood; verify DLQ contents and continuity.
    \item \textbf{Instrument.} Wire the five metrics with alerts and runbook links.
    \item \textbf{Project.} Write the state-growth capacity note with the cardinality sweep.
\end{enumerate}

\subsection*{Deliverables}
\begin{itemize}
    \item Processor and connection definitions, producer/drill scripts, and the metrics evidence.
    \item The two drill reports and the state-growth capacity note.
    \item The runbook entries for lag, state, and DLQ alerts.
\end{itemize}

\subsection*{Acceptance Criteria}
\begin{itemize}
    \item Output parity after restart is byte-exact against the reference run.
    \item The poison flood leaves zero processing gaps and full DLQ accounting.
    \item All five metrics alert correctly under injected conditions.
\end{itemize}

\subsection*{Stretch Goals}
\begin{itemize}
    \item Add a second consumer of the same topic (analytics sink) and prove consumer independence under one side's failure.
    \item Re-implement the velocity window with a hopping window and quantify the state and cost delta.
\end{itemize}
